\documentclass[11pt,a4paper]{article}
% =====================================================================
%  Gemeinsame Vorlage der Arbeitsblätter
%  "Objektorientierte Programmierung mit Java" – Informatik 10 (NTG)
%  Einbinden in jedem Blatt direkt nach \documentclass: % =====================================================================
%  Gemeinsame Vorlage der Arbeitsblätter
%  "Objektorientierte Programmierung mit Java" – Informatik 10 (NTG)
%  Einbinden in jedem Blatt direkt nach \documentclass: % =====================================================================
%  Gemeinsame Vorlage der Arbeitsblätter
%  "Objektorientierte Programmierung mit Java" – Informatik 10 (NTG)
%  Einbinden in jedem Blatt direkt nach \documentclass: \input{ab-vorlage}
%  Übersetzen mit XeLaTeX, LuaLaTeX, pdfLaTeX oder Tectonic.
%
%  Blatt:   tectonic ab-2-5-arrays.tex   (oder xelatex)
%  Lösung:  tectonic ab-2-5-arrays-lsg.tex
%           (Datei enthält nur: \def\mitloesung{}\input{ab-2-5-arrays};
%            füllt alle Lücken rot aus; siehe \lsg, \lsgrest, \lsglinien,
%            \lsgfeld, \lsgkariert, \lsgkarobox, \lsgtext, \lsgoder)
%  Lösungseinträge belegen nie mehr Platz als die Lücke im Blatt.
% =====================================================================
\usepackage[a4paper,left=18mm,right=18mm,top=26mm,bottom=17mm,headheight=15mm,headsep=4mm,footskip=8mm]{geometry}
\usepackage{iftex}
\ifPDFTeX
  \usepackage[T1]{fontenc}
  \usepackage[utf8]{inputenc}
  \usepackage{helvet}
  \usepackage[scaled=0.86]{beramono}
  \renewcommand{\familydefault}{\sfdefault}
\else
  \usepackage{fontspec}
  \setmainfont{texgyreheros}[Extension=.otf,UprightFont=*-regular,BoldFont=*-bold,ItalicFont=*-italic,BoldItalicFont=*-bolditalic]
  \setmonofont{DejaVuSansMono}[Extension=.ttf,UprightFont=*,BoldFont=*-Bold,ItalicFont=*-Oblique,BoldItalicFont=*-BoldOblique,Scale=0.86]
\fi
\usepackage[ngerman,shorthands=off]{babel}
\usepackage{xcolor,graphicx,tabularx,array,enumitem,fancyhdr,tikz,listings,multicol,calc,etoolbox,ragged2e,colortbl,xspace,pifont}
\usepackage[most]{tcolorbox}
\usetikzlibrary{arrows.meta,positioning,calc,decorations.pathreplacing}
\usepackage[hidelinks]{hyperref}
\setlength{\parindent}{0pt}
\setlength{\parskip}{3pt}
\setlist{nosep,leftmargin=*}

% ---------- Lösungsschalter ----------
\newif\ifloesung
\ifdefined\mitloesung\loesungtrue\fi
\newcommand{\lsgfarbe}{\color{red!75!black}}

% ---------- Kopf- und Fußzeile (einheitliche Form) ----------
\newcommand{\lehrkraft}{Hau}
\newcommand{\themenbereich}{2. Objektorientierte Programmierung}
\newcommand{\abnummer}{}
\pagestyle{fancy}\fancyhf{}
\renewcommand{\headrulewidth}{0.4pt}
\fancyhead[L]{\small Klasse 10\\Datum:}
\fancyhead[C]{\small Themenbereich:\\\themenbereich}
\fancyhead[R]{\small \lehrkraft\ifloesung\\{\lsgfarbe\bfseries Lösung}\fi}
\fancyfoot[C]{\scriptsize edu-mrh.de\,·\,Informatik 10\,·\,Blatt \abnummer\,·\,Seite \thepage}
\newcommand{\abtitel}[2]{\renewcommand{\abnummer}{#1}\begin{center}\Large\bfseries #1\quad\underline{#2}\end{center}\vspace{-1mm}}

% ---------- Boxen ----------
\newenvironment{aufgabe}[2][]{\begin{tcolorbox}[enhanced,breakable,colback=white,colframe=black,boxrule=0.7pt,arc=3mm,left=2mm,right=2mm,top=1.5mm,bottom=1.5mm,before skip=2.5mm,after skip=2.5mm]\textbf{Aufgabe #2}\ifstrempty{#1}{}{ \textit{#1}}\par\vspace{0.6mm}}{\end{tcolorbox}}
\newtcolorbox{merke}[1][Merke]{enhanced,breakable,colback=green!5,colframe=green!40!black,boxrule=0.8pt,arc=2mm,left=2mm,right=2mm,top=1.5mm,bottom=1.5mm,before skip=2.5mm,after skip=2.5mm,title={#1},fonttitle=\bfseries,coltitle=white}
\newtcolorbox{hinweis}{enhanced,breakable,colback=yellow!12,colframe=orange!70!black,boxrule=0.6pt,arc=2mm,left=2mm,right=2mm,top=1mm,bottom=1mm,before skip=2mm,after skip=2mm}
\newcommand{\web}[1]{{\small\color{gray}\textit{Website: #1}}}

% ---------- Lücken, Linien, Karos ----------
% Einheitliche Maße für alle Blätter der Sequenz:
%  \linien{n}       n Schreiblinien, Abstand 8 mm, erste Linie 8 mm unter dem Text
%  \kariert{n}      Karoraster 5 mm, n Kästchen hoch, nur ganze Kästchen, mittig
%  \karobox{b}{h}   Karoraster b x h, auf ganze Kästchen gerundet
%  \luecke[l]       Lücke im Fließtext (Standard 4,5 cm, Fachbegriffe mind. 3,5 cm)
%  (jeweils mit Lösungsvariante \lsg…, siehe Kopf der Datei)
%  \restzeile       Lücke bis zum Zeilenende
%  \lueckentext     am Anfang eines Absatzes/Kastens mit Lücken: größerer
%                   Zeilenabstand und Flattersatz, damit man in die Lücken schreiben kann
%  \linienfeld{n}   umrahmtes Feld mit n Schreiblinien (z. B. für Java-Code)
%  \schreibzeile    Strut für Tabellenzellen, in die geschrieben wird (8 mm hoch)
\newlength{\zeilenabstand}\setlength{\zeilenabstand}{8mm}
\newlength{\karo}\setlength{\karo}{5mm}
\tikzset{karo/.style={draw=gray!50,line width=0.3pt},schreiblinie/.style={draw=black!65,line width=0.35pt}}
\newlength{\lsgbreite}
% Lücke: \luecke[Breite] / mit Lösung \lsg[Breite]{Text} (zu langer Text wird verkleinert)
\newcommand{\lsg}[2][4.5cm]{\leavevmode\ifloesung\settowidth{\lsgbreite}{\,#2\,}%
  \ifdim\lsgbreite>#1\relax\rlap{\raisebox{0.8pt}{\resizebox{#1}{!}{\lsgfarbe\,#2\,}}}%
  \else\rlap{\raisebox{0.8pt}{\lsgfarbe\,#2}}\fi\fi\rule[-0.6ex]{#1}{0.4pt}\xspace}
\newcommand{\luecke}[1][4.5cm]{\lsg[#1]{}}
% Lücke bis Zeilenende: \restzeile / \lsgrest{Text}
\newcommand{\lsgrest}[1]{\leavevmode\unskip\hspace{1.5mm}\ifloesung\rlap{\raisebox{0.8pt}{\lsgfarbe\,#1}}\fi\leaders\hbox{\rule[-0.6ex]{0.5pt}{0.4pt}}\hfill\null\par}
\newcommand{\restzeile}{\lsgrest{}}
% nur in der Lösung sichtbar (Platz bleibt frei): \lsgtext{…}; Alternativen: \lsgoder{Lösung}{Blatt}
\newcommand{\lsgtext}[1]{\ifloesung\leavevmode{\lsgfarbe#1}\fi}
\newcommand{\lsgoder}[2]{\ifloesung\leavevmode#1\else#2\fi}
\newcommand{\lueckentext}{\linespread{1.5}\selectfont\raggedright}
\newcommand{\schreibzeile}{\rule[-2.8mm]{0pt}{8mm}}
% Schritt-Nummern wie im Code und auf der Website: \nr{1} = ①
\newcommand{\nr}[1]{\ding{\numexpr171+#1\relax}}
% Schreiblinien: \linien{n} / \lsglinien{n}{Text} (Text liegt auf den Linien)
\newcommand{\lsglinien}[2]{\par\vspace{0.5mm}\noindent\begin{tikzpicture}
  \path (0,0) (\linewidth-0.4pt,-#1\zeilenabstand-1.5mm);
  \foreach \i in {1,...,#1}{\draw[schreiblinie] (0.2pt,-\i\zeilenabstand) -- (\linewidth-0.2pt,-\i\zeilenabstand);}
  \ifloesung\node[overlay,anchor=north west,inner sep=0pt] at (1mm,-\zeilenabstand+1pt+5mm) {\parbox[t]{\linewidth-2mm}{\lsgfarbe\raggedright\setlength{\baselineskip}{\zeilenabstand}\rule{0pt}{5mm}#2}};\fi
  \end{tikzpicture}\par}
\newcommand{\linien}[1]{\lsglinien{#1}{}}
% Codefeld mit Linien: \linienfeld{n} / \lsgfeld{n}{Code, Zeilen mit \\ getrennt}
\newcommand{\lsgfeld}[2]{\noindent\begin{tikzpicture}
  \draw[line width=0.4pt] (0.2pt,0) rectangle (\linewidth-0.2pt,-#1\zeilenabstand-3mm);
  \foreach \i in {1,...,#1}{\draw[schreiblinie] (2mm,-\i\zeilenabstand) -- (\linewidth-2mm,-\i\zeilenabstand);}
  \ifloesung\node[overlay,anchor=north west,inner sep=0pt] at (2.5mm,-\zeilenabstand+1pt+5mm) {\parbox[t]{\linewidth-5mm}{\lsgfarbe\ttfamily\small\setlength{\baselineskip}{\zeilenabstand}\rule{0pt}{5mm}#2}};\fi
  \end{tikzpicture}}
\newcommand{\linienfeld}[1]{\lsgfeld{#1}{}}
% Karoraster; Lösungszeichnung als TikZ-Code, Ursprung oben links, y nach unten negativ:
% \kariert{n} / \lsgkariert{n}{TikZ},  \karobox{b}{h} / \lsgkarobox{b}{h}{TikZ}
\newcommand{\karoraster}[3]{\begin{tikzpicture}\draw[karo] (0,0) grid[step=\karo] (#1\karo,#2\karo);
  \ifloesung\begin{scope}[overlay,shift={(0,#2\karo)},red!75!black,text=red!75!black,line width=0.8pt]#3\end{scope}\fi\end{tikzpicture}}
\newcommand{\lsgkariert}[2]{\par\vspace{1.5mm}\noindent\pgfmathtruncatemacro{\karoX}{floor(\linewidth/\karo)}%
  \makebox[\linewidth]{\karoraster{\karoX}{#1}{#2}}\par\vspace{1mm}}
\newcommand{\kariert}[1]{\lsgkariert{#1}{}}
\newcommand{\lsgkarobox}[3]{\pgfmathtruncatemacro{\karoX}{round((#1)/\karo)}\pgfmathtruncatemacro{\karoY}{round((#2)/\karo)}\karoraster{\karoX}{\karoY}{#3}}
\newcommand{\karobox}[2]{\lsgkarobox{#1}{#2}{}}
% Lösungs-Klassenkarte (für Zeichnungen im Karo): \lsgkarte{Name}{Attribute}{Methoden}
\newcommand{\lsgkarte}[3]{{\lsgfarbe\setlength{\arrayrulewidth}{0.6pt}\arrayrulecolor{red!75!black}\begin{tabular}{|l|}\hline\textbf{#1}\\\hline #2\\\hline #3\\\hline\end{tabular}}}
% Lösungs-Objektkarte (in TikZ als Knoteninhalt): \node[objekt] {\lsgobjekt{name : Klasse}{attr = wert\\…}};
\newcommand{\lsgobjekt}[2]{{\lsgfarbe\small\arrayrulecolor{red!75!black}\begin{tabular}{@{}l@{}}\textbf{#1}\\\hline #2\end{tabular}}}
\tikzset{objekt/.style={draw,rounded corners=3mm,inner sep=2mm,anchor=north west}}
\newcommand{\klassenkarte}[3]{\begin{tabular}{|l|}\hline\textbf{#1}\\\hline #2\\\hline #3\\\hline\end{tabular}}
\newcommand{\code}[1]{\texttt{#1}}

% ---------- Java-Listings ----------
\lstset{language=Java,basicstyle=\ttfamily\small,keywordstyle=\bfseries,commentstyle=\color{gray!80!black},stringstyle=\color{black},showstringspaces=false,columns=flexible,keepspaces=true,tabsize=3,frame=single,framerule=0.4pt,rulecolor=\color{black},xleftmargin=2mm,framexleftmargin=1mm,aboveskip=2mm,belowskip=2mm,breaklines=true,
 literate={ä}{{\"a}}1 {ö}{{\"o}}1 {ü}{{\"u}}1 {Ä}{{\"A}}1 {Ö}{{\"O}}1 {Ü}{{\"U}}1 {ß}{{\ss}}1 {–}{{--}}1 {„}{{\glqq}}1 {“}{{\grqq}}1}
\lstnewenvironment{java}[1][]{\lstset{#1}}{}
% Lückencode: größerer Zeilenabstand, Lücken mit (*\luecke[3cm]*) im Code
\lstnewenvironment{javaluecke}[1][]{\lstset{basicstyle=\linespread{1.5}\small\ttfamily,escapeinside={(*}{*)},#1}}{}

%  Übersetzen mit XeLaTeX, LuaLaTeX, pdfLaTeX oder Tectonic.
%
%  Blatt:   tectonic ab-2-5-arrays.tex   (oder xelatex)
%  Lösung:  tectonic ab-2-5-arrays-lsg.tex
%           (Datei enthält nur: \def\mitloesung{}\documentclass[11pt,a4paper]{article}
\input{ab-vorlage}
\begin{document}
\abtitel{2.5}{Verwaltung gleichartiger Objekte -- Das Array / Feld}

\begin{aufgabe}[Projekt Zufallszahlen]{1}
\textbf{a)} Ergänze die Attribute \code{zahl1} bis \code{zahl3}.\quad \textbf{b)} Weise ihnen im Konstruktor Werte von 1 bis 6 zu.\\
\textbf{c)} Und bei 100 Zahlen? Skizziere, wie ein Speicher aussehen müsste, in dem du „die i-te Zahl“ ansprechen kannst:
\lsgkariert{4}{\begin{scope}[font=\scriptsize]
\foreach \v [count=\i from 0] in {5,2,6,1,4} {\draw (10mm+\i*10mm,-7mm) rectangle ++(10mm,-7mm); \node at (15mm+\i*10mm,-10.5mm) {\v}; \node at (15mm+\i*10mm,-4.5mm) {\i};}
\node at (65mm,-10.5mm) {\dots}; \draw (70mm,-7mm) rectangle ++(10mm,-7mm); \node at (75mm,-10.5mm) {3}; \node at (75mm,-4.5mm) {99};
\node[anchor=west] at (2mm,-4.5mm) {i:}; \node[anchor=west,align=left] at (85mm,-9mm) {ein Name für alle 100 Zahlen,\\Zugriff über die Nummer (Index)};
\end{scope}}
\end{aufgabe}

\begin{minipage}[t]{0.7\linewidth}
\begin{merke}[Felder]\lueckentext
Felder (engl. \lsg[3cm]{Arrays}) dienen der Speicherung einer \lsg[3cm]{festen} \lsg[3cm]{Anzahl} von Werten \lsg[4.5cm]{gleichen Datentyps}, die über einen Index (Positionszahl) adressiert werden können.

Dadurch kann z.\,B. mit einer Wiederholung auf alle Elemente zugegriffen werden, ohne jedes einzeln zu behandeln.
\end{merke}
\end{minipage}\hfill
\begin{minipage}[t]{0.27\linewidth}
\vspace{3mm}
\begin{tabular}{|c|c|}\hline
Index i & \code{textFeld[i]}\\\hline
0 & \code{"Hallo "}\\\hline
1 & \code{"Welt, "}\\\hline
2 & \code{"ein "}\\\hline
3 & \code{"Satz."}\\\hline
\end{tabular}
\end{minipage}

\begin{aufgabe}{2}
Übertrage den Code in das Projekt (TODO 2) und beschreibe die Funktion der Zeilen.
\end{aufgabe}
\begin{javaluecke}
class Zufallszahlen {
   private int[] zahlenFeld;              // (*\lsg[7.5cm]{Deklaration: Feld für int-Werte}*)
   Zufallszahlen() {
      zahlenFeld = new int[8];            // (*\lsg[7.5cm]{Erzeugung: 8 Zellen, alle mit 0 belegt}*)
      zahlenFeld[0] = 5;                  // (*\lsg[7.5cm]{erste Zelle (Index 0) bekommt 5}*)
      zahlenFeld[7] = Random.randint(1, 6); // (*\lsg[7.1cm]{letzte Zelle bekommt Zufallszahl 1--6}*)
   }
}
\end{javaluecke}

\begin{minipage}[t]{0.7\linewidth}
\begin{merke}\lueckentext
\begin{itemize}
\item Felder beginnen ihren Index immer bei \lsg[1.5cm]{0}.
\item Die Größe eines Feldes liefert \code{<Feldname>.length}. Hier also: \lsg[4cm]{\code{zahlenFeld.length}}. Das Array ist \lsg[1.5cm]{8} Zellen groß.
\item Was passiert bei \code{zahlenFeld[8]}?\lsgrest{Laufzeitfehler (Index zu groß)}
\end{itemize}
\end{merke}
\end{minipage}\hfill
\begin{minipage}[t]{0.27\linewidth}
\vspace{3mm}
\begin{tabular}{|c|c|}\hline
Index i & \code{zahlenFeld[i]}\\\hline
0 & 5\\\hline
1 & 0\\\hline
\rule[-1.2mm]{0pt}{5.2mm}2 & \lsgtext{0}\\\hline
\rule[-1.2mm]{0pt}{5.2mm}\lsgtext{3} & \lsgtext{0}\\\hline
\rule[-1.2mm]{0pt}{5.2mm}\lsgtext{4} & \lsgtext{0}\\\hline
\rule[-1.2mm]{0pt}{5.2mm}\lsgtext{5} & \lsgtext{0}\\\hline
\rule[-1.2mm]{0pt}{5.2mm}\lsgtext{6} & \lsgtext{0}\\\hline
\rule[-1.2mm]{0pt}{5.2mm}\lsgtext{7} & \lsgtext{1 bis 6}\\\hline
\end{tabular}
\end{minipage}

\newpage
Fülle die Tabelle auf Seite 1. Warum ist \code{zahlenFeld[7]} die letzte Zelle?
\lsglinien{2}{Der Index beginnt bei 0. Bei 8 Zellen laufen die Indizes von 0 bis 7, der letzte ist also 8 $-$ 1 = 7.}
{\lueckentext Letzte Zelle bei beliebiger Feldgröße:\lsgrest{\code{zahlenFeld[zahlenFeld.length - 1]}}}

\begin{aufgabe}{3}
\textbf{a)} Weise allen Zellen zufällige Werte zu (Wiederholung mit fester Anzahl).\quad
\textbf{b)} 100 Zellen, jede mit der Summe zweier Würfel.\\
Was änderst du an der Schleife aus a), damit sie bei jeder Feldgröße stimmt?
\lsglinien{1}{Bedingung \code{i < zahlenFeld.length} statt einer festen Zahl.}
\end{aufgabe}

\textbf{1:n-Beziehung als Sammlung umsetzen -- Arrays mit Referenzen}

{\centering
\begin{tikzpicture}[x=1cm,y=1cm]
\node[draw,fill=cyan!15,minimum height=6mm] (m) at (0,0) {Mannschaft};
\node[draw,fill=cyan!15,minimum height=6mm] (s) at (5,0) {Spieler};
\draw (m) -- (s) node[midway,above] {\small $<$ spielt in} node[pos=0.08,below] {\small 1} node[pos=0.92,below] {\small n};
\foreach \i/\n in {0/paul,1/tim,2/maja,3/tina} {
  \node[draw,fill=red!10,minimum width=13mm,minimum height=5mm] (z\i) at (6.6+\i*2.2,0.3) {\small\bfseries \i};
  \node[draw,fill=red!10,font=\scriptsize\ttfamily,inner sep=2pt] (o\i) at (6.6+\i*2.2,-1.1) {\n: Spieler};
  \draw[-{Stealth}] (z\i.south) -- (o\i.north);
}
\end{tikzpicture}
\par}
{\lueckentext Auch Referenzen können in Arrays verwaltet werden. Dabei verweist jede Zelle auf ein \lsg[3.5cm]{Objekt} der anderen \lsg[3.5cm]{Klasse}, hier: \lsg[3.5cm]{\code{Spieler}}.\par}

\begin{hinweis}
\lueckentext\textbf{Beachte:} Alle Elemente eines Feldes müssen denselben \lsg[3.5cm]{Datentyp} besitzen, dieser kann aber beliebig gewählt werden: auch \code{boolean}, \code{String}, \code{Rectangle}, \ldots
\end{hinweis}

\begin{merke}[Syntaxübersicht]\small
\begin{tabular}{@{}l l@{}}
\code{<Sichtbarkeit> <Datentyp>[] <Feldname>;} & Deklaration: \code{[]} nach dem Datentyp macht ein Feld daraus.\\
\code{<Feldname> = new <Datentyp>[<Groesse>];} & Erzeugung mit fester Größe (nicht veränderbar).\\
\code{<Feldname>[<Index>]} & Zugriff auf die Zelle an dieser Position.\\
\end{tabular}

\textbf{Beispiele:}\\
lesen: \code{int eineZahl = zahlenFeld[0];}\hfill schreiben: \code{zahlenFeld[0] = eineZahl;}\\
Methode aufrufen: \code{kreisFeld[0].tuEtwas();}\hfill Attribut: \code{kreisFeld[0].xPosition = 100;}
\end{merke}

\begin{minipage}[t]{0.74\linewidth}
\begin{merke}[Mehrdimensionale Arrays]\lueckentext
Arrays können mehrere Dimensionen haben und so eine \lsg[3.5cm]{Tabelle (Raster)} aufspannen.\\
Deklaration und Erzeugung:\lsgrest{\code{Rectangle[][] fuellung}}
\lsgrest{\code{fuellung = new Rectangle[5][5];}}
Zugriff auf Spalte x, Zeile y:\lsgrest{\code{fuellung[x][y]}}
\end{merke}
\end{minipage}\hfill
\begin{minipage}[t]{0.23\linewidth}
\vspace{1mm}\centering
\begin{tikzpicture}[x=6mm,y=6mm,font=\scriptsize]
\ifloesung\fill[red!25] \foreach \x/\y in {1/1,3/1,1/3,2/4,3/3} {(\x,-\y) rectangle (\x+1,-\y-1)};\fi
\draw[step=6mm] (0,0) grid (5,-5);
\foreach \i in {0,...,4} {\node at (\i+0.5,0.4) {\i}; \node at (-0.4,-\i-0.5) {\i};}
\node at (2.5,1.05) {\code{x} $\rightarrow$}; \node[rotate=-90] at (-1.05,-2.5) {\code{y} $\rightarrow$};
\end{tikzpicture}
\end{minipage}

\begin{aufgabe}[Projekt Mosaik]{4}
\textbf{a)} TODO 1 und 2 in \code{Mosaik1}: zweidimensionales Array passender Größe erzeugen und in \code{neuesMosaik()} füllen.\\
\textbf{b)} Auf \code{Mosaik2} wechseln, TODO 3--5 (Smiley). Markiere die Smiley-Zellen \code{fuellung[x][y]} im Raster oben.\lsgtext{\ (z.\,B. [1][1], [3][1], [1][3], [2][4], [3][3])}
\end{aufgabe}
\end{document}
;
%            füllt alle Lücken rot aus; siehe \lsg, \lsgrest, \lsglinien,
%            \lsgfeld, \lsgkariert, \lsgkarobox, \lsgtext, \lsgoder)
%  Lösungseinträge belegen nie mehr Platz als die Lücke im Blatt.
% =====================================================================
\usepackage[a4paper,left=18mm,right=18mm,top=26mm,bottom=17mm,headheight=15mm,headsep=4mm,footskip=8mm]{geometry}
\usepackage{iftex}
\ifPDFTeX
  \usepackage[T1]{fontenc}
  \usepackage[utf8]{inputenc}
  \usepackage{helvet}
  \usepackage[scaled=0.86]{beramono}
  \renewcommand{\familydefault}{\sfdefault}
\else
  \usepackage{fontspec}
  \setmainfont{texgyreheros}[Extension=.otf,UprightFont=*-regular,BoldFont=*-bold,ItalicFont=*-italic,BoldItalicFont=*-bolditalic]
  \setmonofont{DejaVuSansMono}[Extension=.ttf,UprightFont=*,BoldFont=*-Bold,ItalicFont=*-Oblique,BoldItalicFont=*-BoldOblique,Scale=0.86]
\fi
\usepackage[ngerman,shorthands=off]{babel}
\usepackage{xcolor,graphicx,tabularx,array,enumitem,fancyhdr,tikz,listings,multicol,calc,etoolbox,ragged2e,colortbl,xspace,pifont}
\usepackage[most]{tcolorbox}
\usetikzlibrary{arrows.meta,positioning,calc,decorations.pathreplacing}
\usepackage[hidelinks]{hyperref}
\setlength{\parindent}{0pt}
\setlength{\parskip}{3pt}
\setlist{nosep,leftmargin=*}

% ---------- Lösungsschalter ----------
\newif\ifloesung
\ifdefined\mitloesung\loesungtrue\fi
\newcommand{\lsgfarbe}{\color{red!75!black}}

% ---------- Kopf- und Fußzeile (einheitliche Form) ----------
\newcommand{\lehrkraft}{Hau}
\newcommand{\themenbereich}{2. Objektorientierte Programmierung}
\newcommand{\abnummer}{}
\pagestyle{fancy}\fancyhf{}
\renewcommand{\headrulewidth}{0.4pt}
\fancyhead[L]{\small Klasse 10\\Datum:}
\fancyhead[C]{\small Themenbereich:\\\themenbereich}
\fancyhead[R]{\small \lehrkraft\ifloesung\\{\lsgfarbe\bfseries Lösung}\fi}
\fancyfoot[C]{\scriptsize edu-mrh.de\,·\,Informatik 10\,·\,Blatt \abnummer\,·\,Seite \thepage}
\newcommand{\abtitel}[2]{\renewcommand{\abnummer}{#1}\begin{center}\Large\bfseries #1\quad\underline{#2}\end{center}\vspace{-1mm}}

% ---------- Boxen ----------
\newenvironment{aufgabe}[2][]{\begin{tcolorbox}[enhanced,breakable,colback=white,colframe=black,boxrule=0.7pt,arc=3mm,left=2mm,right=2mm,top=1.5mm,bottom=1.5mm,before skip=2.5mm,after skip=2.5mm]\textbf{Aufgabe #2}\ifstrempty{#1}{}{ \textit{#1}}\par\vspace{0.6mm}}{\end{tcolorbox}}
\newtcolorbox{merke}[1][Merke]{enhanced,breakable,colback=green!5,colframe=green!40!black,boxrule=0.8pt,arc=2mm,left=2mm,right=2mm,top=1.5mm,bottom=1.5mm,before skip=2.5mm,after skip=2.5mm,title={#1},fonttitle=\bfseries,coltitle=white}
\newtcolorbox{hinweis}{enhanced,breakable,colback=yellow!12,colframe=orange!70!black,boxrule=0.6pt,arc=2mm,left=2mm,right=2mm,top=1mm,bottom=1mm,before skip=2mm,after skip=2mm}
\newcommand{\web}[1]{{\small\color{gray}\textit{Website: #1}}}

% ---------- Lücken, Linien, Karos ----------
% Einheitliche Maße für alle Blätter der Sequenz:
%  \linien{n}       n Schreiblinien, Abstand 8 mm, erste Linie 8 mm unter dem Text
%  \kariert{n}      Karoraster 5 mm, n Kästchen hoch, nur ganze Kästchen, mittig
%  \karobox{b}{h}   Karoraster b x h, auf ganze Kästchen gerundet
%  \luecke[l]       Lücke im Fließtext (Standard 4,5 cm, Fachbegriffe mind. 3,5 cm)
%  (jeweils mit Lösungsvariante \lsg…, siehe Kopf der Datei)
%  \restzeile       Lücke bis zum Zeilenende
%  \lueckentext     am Anfang eines Absatzes/Kastens mit Lücken: größerer
%                   Zeilenabstand und Flattersatz, damit man in die Lücken schreiben kann
%  \linienfeld{n}   umrahmtes Feld mit n Schreiblinien (z. B. für Java-Code)
%  \schreibzeile    Strut für Tabellenzellen, in die geschrieben wird (8 mm hoch)
\newlength{\zeilenabstand}\setlength{\zeilenabstand}{8mm}
\newlength{\karo}\setlength{\karo}{5mm}
\tikzset{karo/.style={draw=gray!50,line width=0.3pt},schreiblinie/.style={draw=black!65,line width=0.35pt}}
\newlength{\lsgbreite}
% Lücke: \luecke[Breite] / mit Lösung \lsg[Breite]{Text} (zu langer Text wird verkleinert)
\newcommand{\lsg}[2][4.5cm]{\leavevmode\ifloesung\settowidth{\lsgbreite}{\,#2\,}%
  \ifdim\lsgbreite>#1\relax\rlap{\raisebox{0.8pt}{\resizebox{#1}{!}{\lsgfarbe\,#2\,}}}%
  \else\rlap{\raisebox{0.8pt}{\lsgfarbe\,#2}}\fi\fi\rule[-0.6ex]{#1}{0.4pt}\xspace}
\newcommand{\luecke}[1][4.5cm]{\lsg[#1]{}}
% Lücke bis Zeilenende: \restzeile / \lsgrest{Text}
\newcommand{\lsgrest}[1]{\leavevmode\unskip\hspace{1.5mm}\ifloesung\rlap{\raisebox{0.8pt}{\lsgfarbe\,#1}}\fi\leaders\hbox{\rule[-0.6ex]{0.5pt}{0.4pt}}\hfill\null\par}
\newcommand{\restzeile}{\lsgrest{}}
% nur in der Lösung sichtbar (Platz bleibt frei): \lsgtext{…}; Alternativen: \lsgoder{Lösung}{Blatt}
\newcommand{\lsgtext}[1]{\ifloesung\leavevmode{\lsgfarbe#1}\fi}
\newcommand{\lsgoder}[2]{\ifloesung\leavevmode#1\else#2\fi}
\newcommand{\lueckentext}{\linespread{1.5}\selectfont\raggedright}
\newcommand{\schreibzeile}{\rule[-2.8mm]{0pt}{8mm}}
% Schritt-Nummern wie im Code und auf der Website: \nr{1} = ①
\newcommand{\nr}[1]{\ding{\numexpr171+#1\relax}}
% Schreiblinien: \linien{n} / \lsglinien{n}{Text} (Text liegt auf den Linien)
\newcommand{\lsglinien}[2]{\par\vspace{0.5mm}\noindent\begin{tikzpicture}
  \path (0,0) (\linewidth-0.4pt,-#1\zeilenabstand-1.5mm);
  \foreach \i in {1,...,#1}{\draw[schreiblinie] (0.2pt,-\i\zeilenabstand) -- (\linewidth-0.2pt,-\i\zeilenabstand);}
  \ifloesung\node[overlay,anchor=north west,inner sep=0pt] at (1mm,-\zeilenabstand+1pt+5mm) {\parbox[t]{\linewidth-2mm}{\lsgfarbe\raggedright\setlength{\baselineskip}{\zeilenabstand}\rule{0pt}{5mm}#2}};\fi
  \end{tikzpicture}\par}
\newcommand{\linien}[1]{\lsglinien{#1}{}}
% Codefeld mit Linien: \linienfeld{n} / \lsgfeld{n}{Code, Zeilen mit \\ getrennt}
\newcommand{\lsgfeld}[2]{\noindent\begin{tikzpicture}
  \draw[line width=0.4pt] (0.2pt,0) rectangle (\linewidth-0.2pt,-#1\zeilenabstand-3mm);
  \foreach \i in {1,...,#1}{\draw[schreiblinie] (2mm,-\i\zeilenabstand) -- (\linewidth-2mm,-\i\zeilenabstand);}
  \ifloesung\node[overlay,anchor=north west,inner sep=0pt] at (2.5mm,-\zeilenabstand+1pt+5mm) {\parbox[t]{\linewidth-5mm}{\lsgfarbe\ttfamily\small\setlength{\baselineskip}{\zeilenabstand}\rule{0pt}{5mm}#2}};\fi
  \end{tikzpicture}}
\newcommand{\linienfeld}[1]{\lsgfeld{#1}{}}
% Karoraster; Lösungszeichnung als TikZ-Code, Ursprung oben links, y nach unten negativ:
% \kariert{n} / \lsgkariert{n}{TikZ},  \karobox{b}{h} / \lsgkarobox{b}{h}{TikZ}
\newcommand{\karoraster}[3]{\begin{tikzpicture}\draw[karo] (0,0) grid[step=\karo] (#1\karo,#2\karo);
  \ifloesung\begin{scope}[overlay,shift={(0,#2\karo)},red!75!black,text=red!75!black,line width=0.8pt]#3\end{scope}\fi\end{tikzpicture}}
\newcommand{\lsgkariert}[2]{\par\vspace{1.5mm}\noindent\pgfmathtruncatemacro{\karoX}{floor(\linewidth/\karo)}%
  \makebox[\linewidth]{\karoraster{\karoX}{#1}{#2}}\par\vspace{1mm}}
\newcommand{\kariert}[1]{\lsgkariert{#1}{}}
\newcommand{\lsgkarobox}[3]{\pgfmathtruncatemacro{\karoX}{round((#1)/\karo)}\pgfmathtruncatemacro{\karoY}{round((#2)/\karo)}\karoraster{\karoX}{\karoY}{#3}}
\newcommand{\karobox}[2]{\lsgkarobox{#1}{#2}{}}
% Lösungs-Klassenkarte (für Zeichnungen im Karo): \lsgkarte{Name}{Attribute}{Methoden}
\newcommand{\lsgkarte}[3]{{\lsgfarbe\setlength{\arrayrulewidth}{0.6pt}\arrayrulecolor{red!75!black}\begin{tabular}{|l|}\hline\textbf{#1}\\\hline #2\\\hline #3\\\hline\end{tabular}}}
% Lösungs-Objektkarte (in TikZ als Knoteninhalt): \node[objekt] {\lsgobjekt{name : Klasse}{attr = wert\\…}};
\newcommand{\lsgobjekt}[2]{{\lsgfarbe\small\arrayrulecolor{red!75!black}\begin{tabular}{@{}l@{}}\textbf{#1}\\\hline #2\end{tabular}}}
\tikzset{objekt/.style={draw,rounded corners=3mm,inner sep=2mm,anchor=north west}}
\newcommand{\klassenkarte}[3]{\begin{tabular}{|l|}\hline\textbf{#1}\\\hline #2\\\hline #3\\\hline\end{tabular}}
\newcommand{\code}[1]{\texttt{#1}}

% ---------- Java-Listings ----------
\lstset{language=Java,basicstyle=\ttfamily\small,keywordstyle=\bfseries,commentstyle=\color{gray!80!black},stringstyle=\color{black},showstringspaces=false,columns=flexible,keepspaces=true,tabsize=3,frame=single,framerule=0.4pt,rulecolor=\color{black},xleftmargin=2mm,framexleftmargin=1mm,aboveskip=2mm,belowskip=2mm,breaklines=true,
 literate={ä}{{\"a}}1 {ö}{{\"o}}1 {ü}{{\"u}}1 {Ä}{{\"A}}1 {Ö}{{\"O}}1 {Ü}{{\"U}}1 {ß}{{\ss}}1 {–}{{--}}1 {„}{{\glqq}}1 {“}{{\grqq}}1}
\lstnewenvironment{java}[1][]{\lstset{#1}}{}
% Lückencode: größerer Zeilenabstand, Lücken mit (*\luecke[3cm]*) im Code
\lstnewenvironment{javaluecke}[1][]{\lstset{basicstyle=\linespread{1.5}\small\ttfamily,escapeinside={(*}{*)},#1}}{}

%  Übersetzen mit XeLaTeX, LuaLaTeX, pdfLaTeX oder Tectonic.
%
%  Blatt:   tectonic ab-2-5-arrays.tex   (oder xelatex)
%  Lösung:  tectonic ab-2-5-arrays-lsg.tex
%           (Datei enthält nur: \def\mitloesung{}\documentclass[11pt,a4paper]{article}
% =====================================================================
%  Gemeinsame Vorlage der Arbeitsblätter
%  "Objektorientierte Programmierung mit Java" – Informatik 10 (NTG)
%  Einbinden in jedem Blatt direkt nach \documentclass: \input{ab-vorlage}
%  Übersetzen mit XeLaTeX, LuaLaTeX, pdfLaTeX oder Tectonic.
%
%  Blatt:   tectonic ab-2-5-arrays.tex   (oder xelatex)
%  Lösung:  tectonic ab-2-5-arrays-lsg.tex
%           (Datei enthält nur: \def\mitloesung{}\input{ab-2-5-arrays};
%            füllt alle Lücken rot aus; siehe \lsg, \lsgrest, \lsglinien,
%            \lsgfeld, \lsgkariert, \lsgkarobox, \lsgtext, \lsgoder)
%  Lösungseinträge belegen nie mehr Platz als die Lücke im Blatt.
% =====================================================================
\usepackage[a4paper,left=18mm,right=18mm,top=26mm,bottom=17mm,headheight=15mm,headsep=4mm,footskip=8mm]{geometry}
\usepackage{iftex}
\ifPDFTeX
  \usepackage[T1]{fontenc}
  \usepackage[utf8]{inputenc}
  \usepackage{helvet}
  \usepackage[scaled=0.86]{beramono}
  \renewcommand{\familydefault}{\sfdefault}
\else
  \usepackage{fontspec}
  \setmainfont{texgyreheros}[Extension=.otf,UprightFont=*-regular,BoldFont=*-bold,ItalicFont=*-italic,BoldItalicFont=*-bolditalic]
  \setmonofont{DejaVuSansMono}[Extension=.ttf,UprightFont=*,BoldFont=*-Bold,ItalicFont=*-Oblique,BoldItalicFont=*-BoldOblique,Scale=0.86]
\fi
\usepackage[ngerman,shorthands=off]{babel}
\usepackage{xcolor,graphicx,tabularx,array,enumitem,fancyhdr,tikz,listings,multicol,calc,etoolbox,ragged2e,colortbl,xspace,pifont}
\usepackage[most]{tcolorbox}
\usetikzlibrary{arrows.meta,positioning,calc,decorations.pathreplacing}
\usepackage[hidelinks]{hyperref}
\setlength{\parindent}{0pt}
\setlength{\parskip}{3pt}
\setlist{nosep,leftmargin=*}

% ---------- Lösungsschalter ----------
\newif\ifloesung
\ifdefined\mitloesung\loesungtrue\fi
\newcommand{\lsgfarbe}{\color{red!75!black}}

% ---------- Kopf- und Fußzeile (einheitliche Form) ----------
\newcommand{\lehrkraft}{Hau}
\newcommand{\themenbereich}{2. Objektorientierte Programmierung}
\newcommand{\abnummer}{}
\pagestyle{fancy}\fancyhf{}
\renewcommand{\headrulewidth}{0.4pt}
\fancyhead[L]{\small Klasse 10\\Datum:}
\fancyhead[C]{\small Themenbereich:\\\themenbereich}
\fancyhead[R]{\small \lehrkraft\ifloesung\\{\lsgfarbe\bfseries Lösung}\fi}
\fancyfoot[C]{\scriptsize edu-mrh.de\,·\,Informatik 10\,·\,Blatt \abnummer\,·\,Seite \thepage}
\newcommand{\abtitel}[2]{\renewcommand{\abnummer}{#1}\begin{center}\Large\bfseries #1\quad\underline{#2}\end{center}\vspace{-1mm}}

% ---------- Boxen ----------
\newenvironment{aufgabe}[2][]{\begin{tcolorbox}[enhanced,breakable,colback=white,colframe=black,boxrule=0.7pt,arc=3mm,left=2mm,right=2mm,top=1.5mm,bottom=1.5mm,before skip=2.5mm,after skip=2.5mm]\textbf{Aufgabe #2}\ifstrempty{#1}{}{ \textit{#1}}\par\vspace{0.6mm}}{\end{tcolorbox}}
\newtcolorbox{merke}[1][Merke]{enhanced,breakable,colback=green!5,colframe=green!40!black,boxrule=0.8pt,arc=2mm,left=2mm,right=2mm,top=1.5mm,bottom=1.5mm,before skip=2.5mm,after skip=2.5mm,title={#1},fonttitle=\bfseries,coltitle=white}
\newtcolorbox{hinweis}{enhanced,breakable,colback=yellow!12,colframe=orange!70!black,boxrule=0.6pt,arc=2mm,left=2mm,right=2mm,top=1mm,bottom=1mm,before skip=2mm,after skip=2mm}
\newcommand{\web}[1]{{\small\color{gray}\textit{Website: #1}}}

% ---------- Lücken, Linien, Karos ----------
% Einheitliche Maße für alle Blätter der Sequenz:
%  \linien{n}       n Schreiblinien, Abstand 8 mm, erste Linie 8 mm unter dem Text
%  \kariert{n}      Karoraster 5 mm, n Kästchen hoch, nur ganze Kästchen, mittig
%  \karobox{b}{h}   Karoraster b x h, auf ganze Kästchen gerundet
%  \luecke[l]       Lücke im Fließtext (Standard 4,5 cm, Fachbegriffe mind. 3,5 cm)
%  (jeweils mit Lösungsvariante \lsg…, siehe Kopf der Datei)
%  \restzeile       Lücke bis zum Zeilenende
%  \lueckentext     am Anfang eines Absatzes/Kastens mit Lücken: größerer
%                   Zeilenabstand und Flattersatz, damit man in die Lücken schreiben kann
%  \linienfeld{n}   umrahmtes Feld mit n Schreiblinien (z. B. für Java-Code)
%  \schreibzeile    Strut für Tabellenzellen, in die geschrieben wird (8 mm hoch)
\newlength{\zeilenabstand}\setlength{\zeilenabstand}{8mm}
\newlength{\karo}\setlength{\karo}{5mm}
\tikzset{karo/.style={draw=gray!50,line width=0.3pt},schreiblinie/.style={draw=black!65,line width=0.35pt}}
\newlength{\lsgbreite}
% Lücke: \luecke[Breite] / mit Lösung \lsg[Breite]{Text} (zu langer Text wird verkleinert)
\newcommand{\lsg}[2][4.5cm]{\leavevmode\ifloesung\settowidth{\lsgbreite}{\,#2\,}%
  \ifdim\lsgbreite>#1\relax\rlap{\raisebox{0.8pt}{\resizebox{#1}{!}{\lsgfarbe\,#2\,}}}%
  \else\rlap{\raisebox{0.8pt}{\lsgfarbe\,#2}}\fi\fi\rule[-0.6ex]{#1}{0.4pt}\xspace}
\newcommand{\luecke}[1][4.5cm]{\lsg[#1]{}}
% Lücke bis Zeilenende: \restzeile / \lsgrest{Text}
\newcommand{\lsgrest}[1]{\leavevmode\unskip\hspace{1.5mm}\ifloesung\rlap{\raisebox{0.8pt}{\lsgfarbe\,#1}}\fi\leaders\hbox{\rule[-0.6ex]{0.5pt}{0.4pt}}\hfill\null\par}
\newcommand{\restzeile}{\lsgrest{}}
% nur in der Lösung sichtbar (Platz bleibt frei): \lsgtext{…}; Alternativen: \lsgoder{Lösung}{Blatt}
\newcommand{\lsgtext}[1]{\ifloesung\leavevmode{\lsgfarbe#1}\fi}
\newcommand{\lsgoder}[2]{\ifloesung\leavevmode#1\else#2\fi}
\newcommand{\lueckentext}{\linespread{1.5}\selectfont\raggedright}
\newcommand{\schreibzeile}{\rule[-2.8mm]{0pt}{8mm}}
% Schritt-Nummern wie im Code und auf der Website: \nr{1} = ①
\newcommand{\nr}[1]{\ding{\numexpr171+#1\relax}}
% Schreiblinien: \linien{n} / \lsglinien{n}{Text} (Text liegt auf den Linien)
\newcommand{\lsglinien}[2]{\par\vspace{0.5mm}\noindent\begin{tikzpicture}
  \path (0,0) (\linewidth-0.4pt,-#1\zeilenabstand-1.5mm);
  \foreach \i in {1,...,#1}{\draw[schreiblinie] (0.2pt,-\i\zeilenabstand) -- (\linewidth-0.2pt,-\i\zeilenabstand);}
  \ifloesung\node[overlay,anchor=north west,inner sep=0pt] at (1mm,-\zeilenabstand+1pt+5mm) {\parbox[t]{\linewidth-2mm}{\lsgfarbe\raggedright\setlength{\baselineskip}{\zeilenabstand}\rule{0pt}{5mm}#2}};\fi
  \end{tikzpicture}\par}
\newcommand{\linien}[1]{\lsglinien{#1}{}}
% Codefeld mit Linien: \linienfeld{n} / \lsgfeld{n}{Code, Zeilen mit \\ getrennt}
\newcommand{\lsgfeld}[2]{\noindent\begin{tikzpicture}
  \draw[line width=0.4pt] (0.2pt,0) rectangle (\linewidth-0.2pt,-#1\zeilenabstand-3mm);
  \foreach \i in {1,...,#1}{\draw[schreiblinie] (2mm,-\i\zeilenabstand) -- (\linewidth-2mm,-\i\zeilenabstand);}
  \ifloesung\node[overlay,anchor=north west,inner sep=0pt] at (2.5mm,-\zeilenabstand+1pt+5mm) {\parbox[t]{\linewidth-5mm}{\lsgfarbe\ttfamily\small\setlength{\baselineskip}{\zeilenabstand}\rule{0pt}{5mm}#2}};\fi
  \end{tikzpicture}}
\newcommand{\linienfeld}[1]{\lsgfeld{#1}{}}
% Karoraster; Lösungszeichnung als TikZ-Code, Ursprung oben links, y nach unten negativ:
% \kariert{n} / \lsgkariert{n}{TikZ},  \karobox{b}{h} / \lsgkarobox{b}{h}{TikZ}
\newcommand{\karoraster}[3]{\begin{tikzpicture}\draw[karo] (0,0) grid[step=\karo] (#1\karo,#2\karo);
  \ifloesung\begin{scope}[overlay,shift={(0,#2\karo)},red!75!black,text=red!75!black,line width=0.8pt]#3\end{scope}\fi\end{tikzpicture}}
\newcommand{\lsgkariert}[2]{\par\vspace{1.5mm}\noindent\pgfmathtruncatemacro{\karoX}{floor(\linewidth/\karo)}%
  \makebox[\linewidth]{\karoraster{\karoX}{#1}{#2}}\par\vspace{1mm}}
\newcommand{\kariert}[1]{\lsgkariert{#1}{}}
\newcommand{\lsgkarobox}[3]{\pgfmathtruncatemacro{\karoX}{round((#1)/\karo)}\pgfmathtruncatemacro{\karoY}{round((#2)/\karo)}\karoraster{\karoX}{\karoY}{#3}}
\newcommand{\karobox}[2]{\lsgkarobox{#1}{#2}{}}
% Lösungs-Klassenkarte (für Zeichnungen im Karo): \lsgkarte{Name}{Attribute}{Methoden}
\newcommand{\lsgkarte}[3]{{\lsgfarbe\setlength{\arrayrulewidth}{0.6pt}\arrayrulecolor{red!75!black}\begin{tabular}{|l|}\hline\textbf{#1}\\\hline #2\\\hline #3\\\hline\end{tabular}}}
% Lösungs-Objektkarte (in TikZ als Knoteninhalt): \node[objekt] {\lsgobjekt{name : Klasse}{attr = wert\\…}};
\newcommand{\lsgobjekt}[2]{{\lsgfarbe\small\arrayrulecolor{red!75!black}\begin{tabular}{@{}l@{}}\textbf{#1}\\\hline #2\end{tabular}}}
\tikzset{objekt/.style={draw,rounded corners=3mm,inner sep=2mm,anchor=north west}}
\newcommand{\klassenkarte}[3]{\begin{tabular}{|l|}\hline\textbf{#1}\\\hline #2\\\hline #3\\\hline\end{tabular}}
\newcommand{\code}[1]{\texttt{#1}}

% ---------- Java-Listings ----------
\lstset{language=Java,basicstyle=\ttfamily\small,keywordstyle=\bfseries,commentstyle=\color{gray!80!black},stringstyle=\color{black},showstringspaces=false,columns=flexible,keepspaces=true,tabsize=3,frame=single,framerule=0.4pt,rulecolor=\color{black},xleftmargin=2mm,framexleftmargin=1mm,aboveskip=2mm,belowskip=2mm,breaklines=true,
 literate={ä}{{\"a}}1 {ö}{{\"o}}1 {ü}{{\"u}}1 {Ä}{{\"A}}1 {Ö}{{\"O}}1 {Ü}{{\"U}}1 {ß}{{\ss}}1 {–}{{--}}1 {„}{{\glqq}}1 {“}{{\grqq}}1}
\lstnewenvironment{java}[1][]{\lstset{#1}}{}
% Lückencode: größerer Zeilenabstand, Lücken mit (*\luecke[3cm]*) im Code
\lstnewenvironment{javaluecke}[1][]{\lstset{basicstyle=\linespread{1.5}\small\ttfamily,escapeinside={(*}{*)},#1}}{}

\begin{document}
\abtitel{2.5}{Verwaltung gleichartiger Objekte -- Das Array / Feld}

\begin{aufgabe}[Projekt Zufallszahlen]{1}
\textbf{a)} Ergänze die Attribute \code{zahl1} bis \code{zahl3}.\quad \textbf{b)} Weise ihnen im Konstruktor Werte von 1 bis 6 zu.\\
\textbf{c)} Und bei 100 Zahlen? Skizziere, wie ein Speicher aussehen müsste, in dem du „die i-te Zahl“ ansprechen kannst:
\lsgkariert{4}{\begin{scope}[font=\scriptsize]
\foreach \v [count=\i from 0] in {5,2,6,1,4} {\draw (10mm+\i*10mm,-7mm) rectangle ++(10mm,-7mm); \node at (15mm+\i*10mm,-10.5mm) {\v}; \node at (15mm+\i*10mm,-4.5mm) {\i};}
\node at (65mm,-10.5mm) {\dots}; \draw (70mm,-7mm) rectangle ++(10mm,-7mm); \node at (75mm,-10.5mm) {3}; \node at (75mm,-4.5mm) {99};
\node[anchor=west] at (2mm,-4.5mm) {i:}; \node[anchor=west,align=left] at (85mm,-9mm) {ein Name für alle 100 Zahlen,\\Zugriff über die Nummer (Index)};
\end{scope}}
\end{aufgabe}

\begin{minipage}[t]{0.7\linewidth}
\begin{merke}[Felder]\lueckentext
Felder (engl. \lsg[3cm]{Arrays}) dienen der Speicherung einer \lsg[3cm]{festen} \lsg[3cm]{Anzahl} von Werten \lsg[4.5cm]{gleichen Datentyps}, die über einen Index (Positionszahl) adressiert werden können.

Dadurch kann z.\,B. mit einer Wiederholung auf alle Elemente zugegriffen werden, ohne jedes einzeln zu behandeln.
\end{merke}
\end{minipage}\hfill
\begin{minipage}[t]{0.27\linewidth}
\vspace{3mm}
\begin{tabular}{|c|c|}\hline
Index i & \code{textFeld[i]}\\\hline
0 & \code{"Hallo "}\\\hline
1 & \code{"Welt, "}\\\hline
2 & \code{"ein "}\\\hline
3 & \code{"Satz."}\\\hline
\end{tabular}
\end{minipage}

\begin{aufgabe}{2}
Übertrage den Code in das Projekt (TODO 2) und beschreibe die Funktion der Zeilen.
\end{aufgabe}
\begin{javaluecke}
class Zufallszahlen {
   private int[] zahlenFeld;              // (*\lsg[7.5cm]{Deklaration: Feld für int-Werte}*)
   Zufallszahlen() {
      zahlenFeld = new int[8];            // (*\lsg[7.5cm]{Erzeugung: 8 Zellen, alle mit 0 belegt}*)
      zahlenFeld[0] = 5;                  // (*\lsg[7.5cm]{erste Zelle (Index 0) bekommt 5}*)
      zahlenFeld[7] = Random.randint(1, 6); // (*\lsg[7.1cm]{letzte Zelle bekommt Zufallszahl 1--6}*)
   }
}
\end{javaluecke}

\begin{minipage}[t]{0.7\linewidth}
\begin{merke}\lueckentext
\begin{itemize}
\item Felder beginnen ihren Index immer bei \lsg[1.5cm]{0}.
\item Die Größe eines Feldes liefert \code{<Feldname>.length}. Hier also: \lsg[4cm]{\code{zahlenFeld.length}}. Das Array ist \lsg[1.5cm]{8} Zellen groß.
\item Was passiert bei \code{zahlenFeld[8]}?\lsgrest{Laufzeitfehler (Index zu groß)}
\end{itemize}
\end{merke}
\end{minipage}\hfill
\begin{minipage}[t]{0.27\linewidth}
\vspace{3mm}
\begin{tabular}{|c|c|}\hline
Index i & \code{zahlenFeld[i]}\\\hline
0 & 5\\\hline
1 & 0\\\hline
\rule[-1.2mm]{0pt}{5.2mm}2 & \lsgtext{0}\\\hline
\rule[-1.2mm]{0pt}{5.2mm}\lsgtext{3} & \lsgtext{0}\\\hline
\rule[-1.2mm]{0pt}{5.2mm}\lsgtext{4} & \lsgtext{0}\\\hline
\rule[-1.2mm]{0pt}{5.2mm}\lsgtext{5} & \lsgtext{0}\\\hline
\rule[-1.2mm]{0pt}{5.2mm}\lsgtext{6} & \lsgtext{0}\\\hline
\rule[-1.2mm]{0pt}{5.2mm}\lsgtext{7} & \lsgtext{1 bis 6}\\\hline
\end{tabular}
\end{minipage}

\newpage
Fülle die Tabelle auf Seite 1. Warum ist \code{zahlenFeld[7]} die letzte Zelle?
\lsglinien{2}{Der Index beginnt bei 0. Bei 8 Zellen laufen die Indizes von 0 bis 7, der letzte ist also 8 $-$ 1 = 7.}
{\lueckentext Letzte Zelle bei beliebiger Feldgröße:\lsgrest{\code{zahlenFeld[zahlenFeld.length - 1]}}}

\begin{aufgabe}{3}
\textbf{a)} Weise allen Zellen zufällige Werte zu (Wiederholung mit fester Anzahl).\quad
\textbf{b)} 100 Zellen, jede mit der Summe zweier Würfel.\\
Was änderst du an der Schleife aus a), damit sie bei jeder Feldgröße stimmt?
\lsglinien{1}{Bedingung \code{i < zahlenFeld.length} statt einer festen Zahl.}
\end{aufgabe}

\textbf{1:n-Beziehung als Sammlung umsetzen -- Arrays mit Referenzen}

{\centering
\begin{tikzpicture}[x=1cm,y=1cm]
\node[draw,fill=cyan!15,minimum height=6mm] (m) at (0,0) {Mannschaft};
\node[draw,fill=cyan!15,minimum height=6mm] (s) at (5,0) {Spieler};
\draw (m) -- (s) node[midway,above] {\small $<$ spielt in} node[pos=0.08,below] {\small 1} node[pos=0.92,below] {\small n};
\foreach \i/\n in {0/paul,1/tim,2/maja,3/tina} {
  \node[draw,fill=red!10,minimum width=13mm,minimum height=5mm] (z\i) at (6.6+\i*2.2,0.3) {\small\bfseries \i};
  \node[draw,fill=red!10,font=\scriptsize\ttfamily,inner sep=2pt] (o\i) at (6.6+\i*2.2,-1.1) {\n: Spieler};
  \draw[-{Stealth}] (z\i.south) -- (o\i.north);
}
\end{tikzpicture}
\par}
{\lueckentext Auch Referenzen können in Arrays verwaltet werden. Dabei verweist jede Zelle auf ein \lsg[3.5cm]{Objekt} der anderen \lsg[3.5cm]{Klasse}, hier: \lsg[3.5cm]{\code{Spieler}}.\par}

\begin{hinweis}
\lueckentext\textbf{Beachte:} Alle Elemente eines Feldes müssen denselben \lsg[3.5cm]{Datentyp} besitzen, dieser kann aber beliebig gewählt werden: auch \code{boolean}, \code{String}, \code{Rectangle}, \ldots
\end{hinweis}

\begin{merke}[Syntaxübersicht]\small
\begin{tabular}{@{}l l@{}}
\code{<Sichtbarkeit> <Datentyp>[] <Feldname>;} & Deklaration: \code{[]} nach dem Datentyp macht ein Feld daraus.\\
\code{<Feldname> = new <Datentyp>[<Groesse>];} & Erzeugung mit fester Größe (nicht veränderbar).\\
\code{<Feldname>[<Index>]} & Zugriff auf die Zelle an dieser Position.\\
\end{tabular}

\textbf{Beispiele:}\\
lesen: \code{int eineZahl = zahlenFeld[0];}\hfill schreiben: \code{zahlenFeld[0] = eineZahl;}\\
Methode aufrufen: \code{kreisFeld[0].tuEtwas();}\hfill Attribut: \code{kreisFeld[0].xPosition = 100;}
\end{merke}

\begin{minipage}[t]{0.74\linewidth}
\begin{merke}[Mehrdimensionale Arrays]\lueckentext
Arrays können mehrere Dimensionen haben und so eine \lsg[3.5cm]{Tabelle (Raster)} aufspannen.\\
Deklaration und Erzeugung:\lsgrest{\code{Rectangle[][] fuellung}}
\lsgrest{\code{fuellung = new Rectangle[5][5];}}
Zugriff auf Spalte x, Zeile y:\lsgrest{\code{fuellung[x][y]}}
\end{merke}
\end{minipage}\hfill
\begin{minipage}[t]{0.23\linewidth}
\vspace{1mm}\centering
\begin{tikzpicture}[x=6mm,y=6mm,font=\scriptsize]
\ifloesung\fill[red!25] \foreach \x/\y in {1/1,3/1,1/3,2/4,3/3} {(\x,-\y) rectangle (\x+1,-\y-1)};\fi
\draw[step=6mm] (0,0) grid (5,-5);
\foreach \i in {0,...,4} {\node at (\i+0.5,0.4) {\i}; \node at (-0.4,-\i-0.5) {\i};}
\node at (2.5,1.05) {\code{x} $\rightarrow$}; \node[rotate=-90] at (-1.05,-2.5) {\code{y} $\rightarrow$};
\end{tikzpicture}
\end{minipage}

\begin{aufgabe}[Projekt Mosaik]{4}
\textbf{a)} TODO 1 und 2 in \code{Mosaik1}: zweidimensionales Array passender Größe erzeugen und in \code{neuesMosaik()} füllen.\\
\textbf{b)} Auf \code{Mosaik2} wechseln, TODO 3--5 (Smiley). Markiere die Smiley-Zellen \code{fuellung[x][y]} im Raster oben.\lsgtext{\ (z.\,B. [1][1], [3][1], [1][3], [2][4], [3][3])}
\end{aufgabe}
\end{document}
;
%            füllt alle Lücken rot aus; siehe \lsg, \lsgrest, \lsglinien,
%            \lsgfeld, \lsgkariert, \lsgkarobox, \lsgtext, \lsgoder)
%  Lösungseinträge belegen nie mehr Platz als die Lücke im Blatt.
% =====================================================================
\usepackage[a4paper,left=18mm,right=18mm,top=26mm,bottom=17mm,headheight=15mm,headsep=4mm,footskip=8mm]{geometry}
\usepackage{iftex}
\ifPDFTeX
  \usepackage[T1]{fontenc}
  \usepackage[utf8]{inputenc}
  \usepackage{helvet}
  \usepackage[scaled=0.86]{beramono}
  \renewcommand{\familydefault}{\sfdefault}
\else
  \usepackage{fontspec}
  \setmainfont{texgyreheros}[Extension=.otf,UprightFont=*-regular,BoldFont=*-bold,ItalicFont=*-italic,BoldItalicFont=*-bolditalic]
  \setmonofont{DejaVuSansMono}[Extension=.ttf,UprightFont=*,BoldFont=*-Bold,ItalicFont=*-Oblique,BoldItalicFont=*-BoldOblique,Scale=0.86]
\fi
\usepackage[ngerman,shorthands=off]{babel}
\usepackage{xcolor,graphicx,tabularx,array,enumitem,fancyhdr,tikz,listings,multicol,calc,etoolbox,ragged2e,colortbl,xspace,pifont}
\usepackage[most]{tcolorbox}
\usetikzlibrary{arrows.meta,positioning,calc,decorations.pathreplacing}
\usepackage[hidelinks]{hyperref}
\setlength{\parindent}{0pt}
\setlength{\parskip}{3pt}
\setlist{nosep,leftmargin=*}

% ---------- Lösungsschalter ----------
\newif\ifloesung
\ifdefined\mitloesung\loesungtrue\fi
\newcommand{\lsgfarbe}{\color{red!75!black}}

% ---------- Kopf- und Fußzeile (einheitliche Form) ----------
\newcommand{\lehrkraft}{Hau}
\newcommand{\themenbereich}{2. Objektorientierte Programmierung}
\newcommand{\abnummer}{}
\pagestyle{fancy}\fancyhf{}
\renewcommand{\headrulewidth}{0.4pt}
\fancyhead[L]{\small Klasse 10\\Datum:}
\fancyhead[C]{\small Themenbereich:\\\themenbereich}
\fancyhead[R]{\small \lehrkraft\ifloesung\\{\lsgfarbe\bfseries Lösung}\fi}
\fancyfoot[C]{\scriptsize edu-mrh.de\,·\,Informatik 10\,·\,Blatt \abnummer\,·\,Seite \thepage}
\newcommand{\abtitel}[2]{\renewcommand{\abnummer}{#1}\begin{center}\Large\bfseries #1\quad\underline{#2}\end{center}\vspace{-1mm}}

% ---------- Boxen ----------
\newenvironment{aufgabe}[2][]{\begin{tcolorbox}[enhanced,breakable,colback=white,colframe=black,boxrule=0.7pt,arc=3mm,left=2mm,right=2mm,top=1.5mm,bottom=1.5mm,before skip=2.5mm,after skip=2.5mm]\textbf{Aufgabe #2}\ifstrempty{#1}{}{ \textit{#1}}\par\vspace{0.6mm}}{\end{tcolorbox}}
\newtcolorbox{merke}[1][Merke]{enhanced,breakable,colback=green!5,colframe=green!40!black,boxrule=0.8pt,arc=2mm,left=2mm,right=2mm,top=1.5mm,bottom=1.5mm,before skip=2.5mm,after skip=2.5mm,title={#1},fonttitle=\bfseries,coltitle=white}
\newtcolorbox{hinweis}{enhanced,breakable,colback=yellow!12,colframe=orange!70!black,boxrule=0.6pt,arc=2mm,left=2mm,right=2mm,top=1mm,bottom=1mm,before skip=2mm,after skip=2mm}
\newcommand{\web}[1]{{\small\color{gray}\textit{Website: #1}}}

% ---------- Lücken, Linien, Karos ----------
% Einheitliche Maße für alle Blätter der Sequenz:
%  \linien{n}       n Schreiblinien, Abstand 8 mm, erste Linie 8 mm unter dem Text
%  \kariert{n}      Karoraster 5 mm, n Kästchen hoch, nur ganze Kästchen, mittig
%  \karobox{b}{h}   Karoraster b x h, auf ganze Kästchen gerundet
%  \luecke[l]       Lücke im Fließtext (Standard 4,5 cm, Fachbegriffe mind. 3,5 cm)
%  (jeweils mit Lösungsvariante \lsg…, siehe Kopf der Datei)
%  \restzeile       Lücke bis zum Zeilenende
%  \lueckentext     am Anfang eines Absatzes/Kastens mit Lücken: größerer
%                   Zeilenabstand und Flattersatz, damit man in die Lücken schreiben kann
%  \linienfeld{n}   umrahmtes Feld mit n Schreiblinien (z. B. für Java-Code)
%  \schreibzeile    Strut für Tabellenzellen, in die geschrieben wird (8 mm hoch)
\newlength{\zeilenabstand}\setlength{\zeilenabstand}{8mm}
\newlength{\karo}\setlength{\karo}{5mm}
\tikzset{karo/.style={draw=gray!50,line width=0.3pt},schreiblinie/.style={draw=black!65,line width=0.35pt}}
\newlength{\lsgbreite}
% Lücke: \luecke[Breite] / mit Lösung \lsg[Breite]{Text} (zu langer Text wird verkleinert)
\newcommand{\lsg}[2][4.5cm]{\leavevmode\ifloesung\settowidth{\lsgbreite}{\,#2\,}%
  \ifdim\lsgbreite>#1\relax\rlap{\raisebox{0.8pt}{\resizebox{#1}{!}{\lsgfarbe\,#2\,}}}%
  \else\rlap{\raisebox{0.8pt}{\lsgfarbe\,#2}}\fi\fi\rule[-0.6ex]{#1}{0.4pt}\xspace}
\newcommand{\luecke}[1][4.5cm]{\lsg[#1]{}}
% Lücke bis Zeilenende: \restzeile / \lsgrest{Text}
\newcommand{\lsgrest}[1]{\leavevmode\unskip\hspace{1.5mm}\ifloesung\rlap{\raisebox{0.8pt}{\lsgfarbe\,#1}}\fi\leaders\hbox{\rule[-0.6ex]{0.5pt}{0.4pt}}\hfill\null\par}
\newcommand{\restzeile}{\lsgrest{}}
% nur in der Lösung sichtbar (Platz bleibt frei): \lsgtext{…}; Alternativen: \lsgoder{Lösung}{Blatt}
\newcommand{\lsgtext}[1]{\ifloesung\leavevmode{\lsgfarbe#1}\fi}
\newcommand{\lsgoder}[2]{\ifloesung\leavevmode#1\else#2\fi}
\newcommand{\lueckentext}{\linespread{1.5}\selectfont\raggedright}
\newcommand{\schreibzeile}{\rule[-2.8mm]{0pt}{8mm}}
% Schritt-Nummern wie im Code und auf der Website: \nr{1} = ①
\newcommand{\nr}[1]{\ding{\numexpr171+#1\relax}}
% Schreiblinien: \linien{n} / \lsglinien{n}{Text} (Text liegt auf den Linien)
\newcommand{\lsglinien}[2]{\par\vspace{0.5mm}\noindent\begin{tikzpicture}
  \path (0,0) (\linewidth-0.4pt,-#1\zeilenabstand-1.5mm);
  \foreach \i in {1,...,#1}{\draw[schreiblinie] (0.2pt,-\i\zeilenabstand) -- (\linewidth-0.2pt,-\i\zeilenabstand);}
  \ifloesung\node[overlay,anchor=north west,inner sep=0pt] at (1mm,-\zeilenabstand+1pt+5mm) {\parbox[t]{\linewidth-2mm}{\lsgfarbe\raggedright\setlength{\baselineskip}{\zeilenabstand}\rule{0pt}{5mm}#2}};\fi
  \end{tikzpicture}\par}
\newcommand{\linien}[1]{\lsglinien{#1}{}}
% Codefeld mit Linien: \linienfeld{n} / \lsgfeld{n}{Code, Zeilen mit \\ getrennt}
\newcommand{\lsgfeld}[2]{\noindent\begin{tikzpicture}
  \draw[line width=0.4pt] (0.2pt,0) rectangle (\linewidth-0.2pt,-#1\zeilenabstand-3mm);
  \foreach \i in {1,...,#1}{\draw[schreiblinie] (2mm,-\i\zeilenabstand) -- (\linewidth-2mm,-\i\zeilenabstand);}
  \ifloesung\node[overlay,anchor=north west,inner sep=0pt] at (2.5mm,-\zeilenabstand+1pt+5mm) {\parbox[t]{\linewidth-5mm}{\lsgfarbe\ttfamily\small\setlength{\baselineskip}{\zeilenabstand}\rule{0pt}{5mm}#2}};\fi
  \end{tikzpicture}}
\newcommand{\linienfeld}[1]{\lsgfeld{#1}{}}
% Karoraster; Lösungszeichnung als TikZ-Code, Ursprung oben links, y nach unten negativ:
% \kariert{n} / \lsgkariert{n}{TikZ},  \karobox{b}{h} / \lsgkarobox{b}{h}{TikZ}
\newcommand{\karoraster}[3]{\begin{tikzpicture}\draw[karo] (0,0) grid[step=\karo] (#1\karo,#2\karo);
  \ifloesung\begin{scope}[overlay,shift={(0,#2\karo)},red!75!black,text=red!75!black,line width=0.8pt]#3\end{scope}\fi\end{tikzpicture}}
\newcommand{\lsgkariert}[2]{\par\vspace{1.5mm}\noindent\pgfmathtruncatemacro{\karoX}{floor(\linewidth/\karo)}%
  \makebox[\linewidth]{\karoraster{\karoX}{#1}{#2}}\par\vspace{1mm}}
\newcommand{\kariert}[1]{\lsgkariert{#1}{}}
\newcommand{\lsgkarobox}[3]{\pgfmathtruncatemacro{\karoX}{round((#1)/\karo)}\pgfmathtruncatemacro{\karoY}{round((#2)/\karo)}\karoraster{\karoX}{\karoY}{#3}}
\newcommand{\karobox}[2]{\lsgkarobox{#1}{#2}{}}
% Lösungs-Klassenkarte (für Zeichnungen im Karo): \lsgkarte{Name}{Attribute}{Methoden}
\newcommand{\lsgkarte}[3]{{\lsgfarbe\setlength{\arrayrulewidth}{0.6pt}\arrayrulecolor{red!75!black}\begin{tabular}{|l|}\hline\textbf{#1}\\\hline #2\\\hline #3\\\hline\end{tabular}}}
% Lösungs-Objektkarte (in TikZ als Knoteninhalt): \node[objekt] {\lsgobjekt{name : Klasse}{attr = wert\\…}};
\newcommand{\lsgobjekt}[2]{{\lsgfarbe\small\arrayrulecolor{red!75!black}\begin{tabular}{@{}l@{}}\textbf{#1}\\\hline #2\end{tabular}}}
\tikzset{objekt/.style={draw,rounded corners=3mm,inner sep=2mm,anchor=north west}}
\newcommand{\klassenkarte}[3]{\begin{tabular}{|l|}\hline\textbf{#1}\\\hline #2\\\hline #3\\\hline\end{tabular}}
\newcommand{\code}[1]{\texttt{#1}}

% ---------- Java-Listings ----------
\lstset{language=Java,basicstyle=\ttfamily\small,keywordstyle=\bfseries,commentstyle=\color{gray!80!black},stringstyle=\color{black},showstringspaces=false,columns=flexible,keepspaces=true,tabsize=3,frame=single,framerule=0.4pt,rulecolor=\color{black},xleftmargin=2mm,framexleftmargin=1mm,aboveskip=2mm,belowskip=2mm,breaklines=true,
 literate={ä}{{\"a}}1 {ö}{{\"o}}1 {ü}{{\"u}}1 {Ä}{{\"A}}1 {Ö}{{\"O}}1 {Ü}{{\"U}}1 {ß}{{\ss}}1 {–}{{--}}1 {„}{{\glqq}}1 {“}{{\grqq}}1}
\lstnewenvironment{java}[1][]{\lstset{#1}}{}
% Lückencode: größerer Zeilenabstand, Lücken mit (*\luecke[3cm]*) im Code
\lstnewenvironment{javaluecke}[1][]{\lstset{basicstyle=\linespread{1.5}\small\ttfamily,escapeinside={(*}{*)},#1}}{}

\begin{document}
\abtitel{2.6}{Arrays durchsuchen}

Im Rahmen des Schul-Fußballturniers sollen Statistiken erstellt werden. Janina ergänzt daher für jeden Spieler bzw. jede Spielerin die Anzahl der bisher erzielten Tore. Die Personen einer Mannschaft sind wie bisher in der \code{spielerliste} gespeichert.

\begin{center}
\begin{tikzpicture}[x=1cm,y=1cm]
\foreach \i/\n/\t/\z in {0/Yusuf/4/45,1/Marie/6/90,2/Luca/0/90,3/Tom/0/45} {
  \node[draw=red!60,fill=red!8,minimum width=34mm,minimum height=13mm,anchor=west] at (\i*3.6,0) {};
  \node[fill=red!70,text=white,font=\bfseries,minimum size=5mm,anchor=north west] at (\i*3.6+0.05,0.6) {\i};
  \node[anchor=north west,font=\ttfamily\scriptsize,align=left,inner sep=1pt] at (\i*3.6+0.6,0.62) {\n\\ Tore: \t\\ Min.: \z};
}
\end{tikzpicture}
\end{center}

Beschreibe in Worten (noch kein Java), wie mithilfe der \code{spielerliste} berechnet werden kann, wie viele Tore die Mannschaft insgesamt geschossen hat. Die Nummern findest du später im Code wieder.

{\lueckentext
\nr{1} Vorher:\lsgrest{eine Hilfsvariable \code{summe} mit 0 anlegen}
\nr{2} Wiederhole für\lsgrest{jeden Spieler der \code{spielerliste} (Index 0 bis Länge $-$ 1)}
\nr{3} In jedem Durchlauf:\lsgrest{die Tore dieses Spielers zu \code{summe} addieren}
\nr{4} Am Ende:\lsgrest{\code{summe} zurückgeben}\par}
Überlegt zu zweit: Welche weiteren Tor-Statistiken könnten aus der \code{spielerliste} erstellt werden?
\lsglinien{2}{z.\,B. beste Torschützin / bester Torschütze, Tore pro Spielminute, Durchschnitt pro Spieler, Anzahl der Spieler ohne Tor}

\begin{aufgabe}[Projekt Tore]{1}
\textbf{a)} Ergänze in \code{Mannschaft} die Methode \code{gibToreGesamt()} und gib das Ergebnis im Hauptprogramm aus. Warum geht \code{spielerliste[i].tore} nicht, und was nimmst du stattdessen?
\lsglinien{1}{\code{tore} ist \code{private} (Station 2.4). Stattdessen den Getter: \code{spielerliste[i].gibTore()}}
\end{aufgabe}

\begin{merke}[Bedingungen an Elemente im Array stellen]\lueckentext
Die Kombination von \lsg[4.5cm]{Wiederholung (Schleife)} und \lsg[4.5cm]{bedingter Anweisung} kann genutzt werden, um Bedingungen an die Elemente eines Arrays zu stellen: Eine Anweisung wird nur ausgeführt, wenn die Bedingung für das aktuelle Element erfüllt ist. So lassen sich etwa nur die bisher torlosen Spieler ausgeben.
\end{merke}

\begin{aufgabe}[Projekt Tore]{1}
\textbf{b)} Fülle die Lücken und übertrage den Code ins Hauptprogramm (TODO 2).
\end{aufgabe}
\begin{javaluecke}
for (int i = (*\lsg[1.5cm]{0}*); i < (*\lsg[5cm]{spielerliste.length}*); i(*\lsg[1.5cm]{++}*)) {
   if (spielerliste[i].(*\lsg[4.5cm]{gibTore()}*) == 0) {
      println((*\lsg[6cm]{spielerliste[i]}*));
   }
}
\end{javaluecke}

\newpage
\begin{aufgabe}[Projekt Sportarten]{2}
\textbf{a)} Ergänze in \code{Spieler} die Methode \code{ersetzeSportartTennis()} und teste sie im Hauptprogramm.\\
\textbf{b)} Ergänze \code{spieltSportart(String sportart)} mit sinnvollem Rückgabetyp: \lsg[3cm]{\code{boolean}}. Zeichne zuerst das Struktogramm:
\end{aufgabe}
\lsgkariert{18}{\begin{scope}[shift={(10mm,-12mm)},font=\small]
\node[anchor=south west] at (0,1mm) {\code{boolean spieltSportart(String sportart)}};
\draw (0,0) rectangle (130mm,-62mm);
\draw (0,-8mm) -- (130mm,-8mm); \node[anchor=west] at (2mm,-4mm) {\code{ergebnis = false}};
\draw (8mm,-54mm) -- (8mm,-16mm) -- (130mm,-16mm); \node[anchor=west] at (2mm,-12mm) {Wiederhole für \code{i} von 0 bis \code{sportarten.length - 1}};
\draw (8mm,-16mm) -- (69mm,-30mm) -- (130mm,-16mm); \draw (8mm,-30mm) -- (130mm,-30mm); \draw (69mm,-30mm) -- (69mm,-54mm);
\node at (69mm,-21mm) {\code{sportarten[i].equals(sportart)}}; \node at (13mm,-27.5mm) {wahr}; \node at (125mm,-27.5mm) {falsch};
\node at (38.5mm,-42mm) {\code{ergebnis = true}}; \draw (69mm,-30mm) -- (130mm,-54mm);
\draw (0,-54mm) -- (130mm,-54mm); \node[anchor=west] at (2mm,-58mm) {gib \code{ergebnis} zurück};
\end{scope}}
\end{document}
